\chapter{Konzeption}

In diesem Kapitel wird die Konzeption der mobilen Anwendung von 
\textit{petappoint} beschrieben. Aufbauend auf den in Kapitel~3 
erarbeiteten Anforderungen wird zunächst die Softwarearchitektur 
erläutert, gefolgt vom UI/UX-Konzept sowie der Datenkommunikation zwischen der mobilen Anwendung und dem bestehenden Backend.

\section{Softwarearchitektur}

Softwarearchitektur bezeichnet die organisatorische Struktur eines Softwaresystems, das aus Computerprogrammen, Prozeduren sowie zugehöriger Dokumentation und Daten besteht. Sie beschreibt, wie die einzelnen Komponenten, Module und Teilprogramme eines Systems aufgebaut und miteinander verbunden sind~\cite{ieee1990}.

\subsection{Architekturübersicht}

Die Architektur von \textit{petappoint} folgt dem Client-Server-Prinzip und basiert auf einer Schichtenarchitektur. Diese gliedert sich in drei Schichten: 

\begin{itemize}
    \item \textbf{Präsentationsschicht (Frontend):} Zuständig für die Darstellung der Benutzeroberfläche.
    \item \textbf{Anwendungsschicht (Backend):} Verarbeitet eingehende HTTP-Anfragen, validiert Eingabedaten und steuert den Datenbankzugriff.
    \item \textbf{Datenschicht (Datenbank):} Verantwortlich für die Speicherung und Verwaltung aller persistenten Daten der Anwendung.
\end{itemize}

Die mobile Anwendung wird in der Präsentationsschicht ergänzt und ist somit ausschließlich für die Darstellung und Nutzerinteraktion zuständig. Sie kommuniziert über eine REST-API mit der Geschäftslogikschicht des Backends.

\subsection{Bestehende Backend-Architektur als Grundlage}

TODO

\subsection{Komponentendiagramm}

\subsection{Klassendiagramm}

\subsection{Navigationskonzeptdiagramm}

\subsection{Mockups}